\documentclass[11pt]{article}
\usepackage[margin=1in]{geometry}
\usepackage{amsmath,amssymb,amsthm,microtype}
\usepackage[hidelinks]{hyperref}
\hypersetup{pdftitle={Partial results on primitive binary pentanomials}}
\newtheorem{theorem}{Theorem}
\newtheorem{lemma}[theorem]{Lemma}
\newtheorem{corollary}[theorem]{Corollary}
\newcommand{\F}{\mathbb F_2}
\newcommand{\ord}{\operatorname{ord}}
\title{Partial results on primitive binary pentanomials}
\author{}
\date{}
\begin{document}
\maketitle
\vspace{-2em}
The question is whether every positive degree \(n\) admits a primitive polynomial
in \(\F[X]\) with at most five nonzero coefficients. A primitive polynomial is
monic and irreducible, and its root generates the group of order \(2^n-1\).
For \(n>1\) its constant coefficient is 1, and evaluating at 1 shows that its
weight is odd. Thus a witness in these degrees must be a trinomial or a
pentanomial. Degree one has the witness \(X+1\). The all-degree existence
question remains unresolved here. This exact assertion is Conjecture 4.3.5
in Stephen D. Cohen's section
\href{https://archive.ymsc.tsinghua.edu.cn/pacm_download/672/12637-dingjt-p2.pdf\#page=133}
{\emph{Weights of primitive polynomials}}, which cites Golomb and
Mullen--Shparlinski. The weaker irreducibility question is Conjecture 2.2.5 in the
\href{https://people.math.carleton.ca/~dthomson/Research/pdfs/DThomson-HFFSect22.pdf}
{chapter by David Thomson in the \emph{Handbook of Finite Fields}}.

\section{A family with two free exponents}
Put
\[
 F_{n,u,v}(X)=X^n+(X^u+1)(X^v+1)
             =X^n+X^{u+v}+X^u+X^v+1 .
\]
Throughout this section \(u,v\) are distinct positive integers with \(u+v<n\).

\begin{theorem}
For every \(5\le n\le256\), some \(F_{n,u,v}\) is primitive.
Consequently every \(1\le n\le256\) has a primitive polynomial of weight at most five.
\end{theorem}
\begin{proof}
The witness tables are imported by \texttt{GF2Sparse/Coverage.lean}.
For example, \(F_{24,1,3}=X^{24}+X^4+X^3+X+1\) and
\(F_{256,2,27}=X^{256}+X^{29}+X^{27}+X^2+1\).
For each candidate \(f\), write \(r=X\bmod f\) and \(N=2^n-1\).
The certificate checks \(r^N=1\) and \(r^{N/p}\ne1\) for every prime divisor
\(p\) of \(N\). The first equality makes \(r\) a unit of order dividing \(N\).
Every proper divisor of \(N\) divides \(N/p\) for some such \(p\).
Thus the \(N\) distinct powers of \(r\) exhaust the nonzero elements of the
\(2^n\)-element quotient. The quotient is a field and \(f\) is primitive.
The Lean proof also checks each integer factorization and the primality of its
factors. Direct certificates cover degrees 1 through 4.
\end{proof}
The public \href{https://people.math.carleton.ca/~daniel/hff/}{handbook tables}
give minimum-weight primitive polynomials for degrees 2 through 640.
The relation \(u+v\) between the interior exponents also appears in earlier
work. Table 4.8 of the \emph{Handbook of Applied Cryptography} uses it for
its pentanomial entries, giving a trinomial when one is available.
Theorem 1 verifies a pentanomial of this shape throughout \(5\le n\le256\),
including degrees admitting a primitive trinomial.
\href{https://cacr.uwaterloo.ca/hac/about/chap4.pdf\#page=37}{Menezes, van Oorschot and Vanstone, chapter 4}.

\begin{theorem}\label{thm:sieve}
Let \(B\ge1\) and \(L=\mathop{\rm lcm}_{1\le d\le B}(2^d-1)\).
For every \(n>L+2\), \(F_{n,L,2}\) has exactly five terms, is squarefree,
and all its irreducible factors have degree greater than \(B\).
\end{theorem}
\begin{proof}
The integer \(L\) is odd, so \(L\ne2\). If a nonzero \(\alpha\) has degree
\(d\le B\) over \(\F\), then \(\alpha^{2^d-1}=1\). Hence
\(F_{n,L,2}(\alpha)=\alpha^n\ne0\). The constant coefficient excludes zero.

More generally \(F_{n,u,v}\) is squarefree when \(u\) is odd and \(v\) is even.
At a root \(r\), both \(r\) and \(s=r^v+1\) are nonzero.
For even \(n\), the derivative at \(r\) equals \(r^{u-1}s\).
For odd \(n\), the identity \(F(r)+rF'(r)=s\) gives \(F'(r)\ne0\).
\end{proof}

\begin{theorem}\label{thm:order}
For \(d=u+v\), the order of \(X\) in the units of \(\F[X]/F_{n,u,v}\) is at least
\[
 C(n,d)=\sum_{j=0}^{\lfloor(n-1)/d\rfloor}(n-dj)\ \ge\ \frac{n^2}{2d}.
\]
\end{theorem}
\begin{proof}
Put \(g=(X^u+1)(X^v+1)\). Since \(F(0)=1\), \(X\) is a unit in the quotient.
The residue of \(X^{i+nj}\) is represented by \(X^ig^j\).
If \(i+dj<n\), this representative has degree less than \(n\).
Its least exponent is \(i\) and its largest is \(i+dj\), so it determines
both \(i\) and \(j\). Distinct representatives of degree less than \(n\)
give distinct residues. Counting the pairs gives the sum. It bounds below the
area under \(\max(n-dx,0)\) for \(x\ge0\).
\end{proof}
In particular Theorem~\ref{thm:sieve} gives order at least
\(n^2/(2(L+2))\) for every \(n>L+2\). The target order is \(2^n-1\).

\paragraph{Factor parity.}
Take \(B\ge3\). For all sufficiently large \(n\equiv1,4,6,7\bmod8\),
\(F_{n,1,4L}\) is squarefree, avoids factors of degree at most \(B\),
and has an \emph{odd} number of irreducible factors. The first two claims
follow from the derivative and \(r^{4L}=1\) in the small fields. For parity,
use Theorems 3.1 and 3.2 of Kim, Pak and Sin with even exponent \(4L\),
divisible by four. Their size condition holds for \(n\ge8L+2\).
For \(n\) a positive multiple of \(8L\), the same three conclusions hold
for \(F_{n,1,2L}\). Apply Theorem 3.2 with leading degree \(m=n\)
and even exponent \(h=2L\). Here \(L\ge21\),
\(h\le\lfloor n/2\rfloor-1\), \(h\equiv2\bmod4\), \(h\ne6\), and
\((2n+4)\bmod h=4\). None of the even-factor cases applies.
Reducibility of this candidate would therefore require at least three factors,
each of degree greater than \(B\). This deduction uses the published parity
theorem. \href{https://arxiv.org/abs/1406.7597}{Kim, Pak and Sin, \emph{Swan-like reducibility for Type I pentanomials}}.

\section{How many pentanomials survive a finite sieve?}
\begin{theorem}\label{thm:many}
Fix \(B\) and write \(L\) as in Theorem~\ref{thm:sieve}.
For \(n>6LM\) there are at least \(LM^3\) distinct squarefree degree-\(n\)
pentanomials whose irreducible factors all have degree greater than \(B\).
In particular their number is \(\Omega_B(n^3)\) in every large degree.
\end{theorem}
\begin{proof}
Use \(X^n+X^c+X^b+X^a+1\). Let \(\delta=1\) for even \(n\) and
\(\delta=2\) for odd \(n\). Choose integers \(0\le i<M,\ 1\le j\le LM,\
1\le k\le M\), and put
\[
  a=L(2i+\delta),\qquad b=2LM+2j,\qquad c=b+2Lk.
\]
Then \(0<a\le 2LM<b<c\le 6LM<n\), so the support determines \(a,b,c\),
which in turn determine \(i,j,k\). Exactly one of \(n,a,b,c\) is odd.
The derivative is a monomial, and the constant term is 1, proving
squarefreeness. At every nonzero \(\alpha\) of degree at most \(B\),
\(\alpha^a=1\) and \(\alpha^c=\alpha^b\), so the value is \(\alpha^n\ne0\).
Take \(M=\lfloor(n-1)/(6L)\rfloor\) for the asymptotic claim.
\end{proof}
There are \(\binom{n-1}{3}\) degree-\(n\) pentanomials with constant
coefficient 1. The lower limit of the fraction of these passing the
squarefreeness and small-factor tests is at least \(1/(36L^2)\).
The cutoff can grow with \(n\). For \(0<\varepsilon<1\), take the largest
integer \(B\) with \(B(B+1)\le\varepsilon\log_2 n\). Then
\(L^2\le\prod_{j=1}^B(2^j-1)^2\le n^\varepsilon\).
In every sufficiently large degree there are therefore
\(\gg n^{3-\varepsilon}\) squarefree pentanomials with no irreducible
factor of degree at most this \(B\), which grows as
\(\sqrt{\varepsilon\log_2 n}\).

\section{Passing the factor-parity test in every large degree}
\begin{theorem}\label{thm:parity}
Fix \(B\ge1\), put \(L=\mathop{\rm lcm}_{1\le d\le B}(2^d-1)\), and let
\(n>128L\). There are at least
\[
                      \frac{n^2}{6144L^2}
\]
distinct squarefree degree-\(n\) binary pentanomials with constant term 1,
an odd number of irreducible factors, and no factor of degree from 1 to \(B\).
The positive exponents of each have greatest common divisor 1, and the
support is a Sidon set. Its fifteen unordered pair sums, including doubles,
are all distinct.
\end{theorem}
\begin{proof}
Every nonzero element of degree at most \(B\) has \(L\)-th power 1.
In both constructions below two pairs of terms cancel at every such
element, leaving the leading monomial. Each polynomial has exactly one
odd exponent, so its derivative is a monomial. Its constant term 1
then gives squarefreeness.

\emph{Odd \(n\).} Choose a multiple \(h\) of \(L\) in \((n/4,n/3)\).
Require \(h\) to be odd and choose its residue modulo 4 so that
\[
        4\mid n-3h \quad\Longleftrightarrow\quad n\equiv3,5\pmod8 .
\]
This selects one residue of \(h\) modulo \(4L\), since \(L\) is odd.
Choose \(3n/4<c<n\) with \(c\equiv8\bmod 8L\), and take the integer lift
\[
              F(X)=1+X^8+X^{2h}+X^c+X^n .
\]
The cancelling pairs are \((0,2h)\) and \((8,c)\). The exponents \(8,n\)
are coprime.

Write \(\beta_i\) for the inverses of the complex roots of \(F\), with
multiplicity, and \(P_j=\sum_i\beta_i^j\). They are algebraic units.
Set \(k=n-2h\), \(w=n-3h\), so that \(k\) is odd, \(2k<n\), and
\(0<w<h<c/2\). The product formula for the discriminant, with the
derivative taken modulo 8, gives
\[
 \operatorname{disc}(F)\equiv(-1)^{n(n-1)/2}n^n
 \left(1+\frac{2h}{n}P_k+
       \left(\frac{2h}{n}\right)^2\frac{P_k^2-P_{2k}}2\right)\pmod8 .
\]
Here \(n\) is invertible modulo 8, and the half in the last term is the
integer second elementary symmetric function of the \(\beta_i^k\).
Indeed only \(nX^{n-1}\) and \(2hX^{2h-1}\) survive in \(F'\) modulo 8,
and the product of all \((n-1)\)-st powers of the roots is 1.

Put \(G(Y)=1+Y^4+Y^h+Y^{c/2}\). From
\(\sum_{j\ge1}P_jX^j=-XF'(X)/F(X)\) and
\(F(X)\equiv G(X^2)\pmod{X^n}\), one gets \(P_k=0\) and
\[
 \frac{P_{2k}}2=[Y^k]\frac{-YG'(Y)}{G(Y)}
 \equiv [Y^{w}]\frac1{1+Y^4}=\mathbf 1_{4\mid w}\pmod2 .
\]
In the last step \(YG'\equiv Y^h\bmod2\), and \(w<h,c/2\).
Thus the discriminant is \((-1)^{n(n-1)/2}n^n(1+4\mathbf 1_{4\mid w})\)
modulo 8. Its first factor is 5 for \(n\equiv3,5\bmod8\) and 1 for
\(n\equiv1,7\bmod8\). Our choice makes the discriminant 1 in all four
cases. The Stickelberger--Swan theorem gives odd factor count for the
binary reduction since \(n\) is odd.
\href{https://doi.org/10.2140/pjm.1962.12.1099}{Swan, Theorem 1}.

For the Sidon property, \(n\) is the sole odd exponent and is largest.
An equality of two different unordered pair sums must therefore use
only \(0,8,2h,c\). Among these four, \(2h>16\) and \(c>2h+8\).
The only remaining possibilities are \(c=4h\) or \(c=4h-8\).
The first exceeds \(n\) and the second is 4 modulo 8.

\emph{Even \(n\).} Take \(a=4\) if \(n\equiv4\bmod8\), and \(a=2\)
otherwise. Choose \(n/16<b<n/8\) with \(L\mid b\). Require
\(b\equiv n/2\bmod4\) if \(n\equiv2\bmod4\), and \(b\equiv1\bmod4\)
otherwise. Choose \(n/4<c<n/3\) with \(c\equiv a\bmod 2L\), and use
\[
                    1+X^a+X^b+X^c+X^n .
\]
The pairs \((0,b),(a,c)\) cancel. The odd exponent \(b\) is coprime
to \(a\), a power of two. Set \(k=b-a,\ t=n-b,\ v=c-b\), and take the
integer lift \(F\) with all five coefficients positive. For its complex
roots \(\alpha_i\) put \(S_j=\sum_i\alpha_i^j\), also for negative \(j\).
The needed power sums are
\[
\begin{gathered}
 S_{-k}=0,\qquad S_t=-t,\qquad S_{2t}=t,\\
 S_v=S_{2v}=S_{t-k}=S_{v-k}=S_{t+v}=0,\qquad
 S_{-2k}=-2\quad\text{if }a=2 .
\end{gathered}
\]
For positive indices use \(-X(F^*)'/F^*\), where
\(F^*=1+X^{n-c}+X^t+X^{n-a}+X^n\).
Since \(3(n-c)>2n\), at most two positive exponents of \(F^*\) can
contribute to \(S_{2t}\). The only pair summing to \(2t\) is \((t,t)\),
because \(2a<2b<c\). The indices \(v,2v,v-k\) lie strictly between
0 and \(n-c\), \(t-k\) lies between \(n-c\) and \(t\), and
\(n<t+v<2(n-c)\). This gives every listed positive-index value.
For \(S_{-k}\), use \(-XF'/F\). The index \(k<b\) is odd and all exponents below
\(b\) are even. When \(a=2\), the coefficient at \(2k<2b<c\) comes
entirely from \(-2X^2/(1+X^2)\). A single use of \(X^b\) has odd
degree and two uses would exceed \(2k\).

Factor \(bX^{b-1}\) from the derivative evaluated at each root. This
leaves \(1+(a/b)\alpha_i^{-k}+(c/b)\alpha_i^v+(n/b)\alpha_i^t\).
All three coefficients are even modulo 8. In the product over \(i\)
retain terms of total degree at most two. The displayed power sums
make all mixed terms zero, using
\(\sum_{i\ne j}\alpha_i^r\alpha_j^s=S_rS_s-S_{r+s}\).
The terms in \(n/b\) give \(1-t(n/b)+\binom t2(n/b)^2
\equiv(1-n/b)^t\bmod8\). The terms in \(c/b\) vanish, and the
quadratic term in \(a/b\) gives \(4\mathbf1_{a=2}\).
Since \(b^n\equiv1\bmod8\) and \(t\) is odd, this yields
\[
 \operatorname{disc}(F)\equiv
 (-1)^{n/2}\bigl(1-n/b+4\mathbf1_{a=2}\bigr)\pmod8 .
\]
For \(n\equiv0,4\bmod8\) the respective choices \(a=2,4\) both give 5.
For \(n\equiv2\bmod4\) we have \(n/b\equiv2\bmod8\) and \(a=2\),
again giving 5. Stickelberger--Swan now gives an odd number of factors.
The same parity conclusion follows independently from cases (1c) and
(2d) of Theorem 6 of Ryul Kim and Yun Mi Kim, using \(c>2b\) and
\(n>3c-2b\).
\href{https://doi.org/10.1155/2014/482837}{Kim and Kim,
\emph{The Reducibility of a Special Binary Pentanomial}}.

Here \(2a<b,\ 2b<c,\ 2c<n\). At each step the next exponent exceeds
every sum of two earlier exponents. The new sums are distinct from
each other as well, which proves the Sidon assertion.

\emph{Counting.} An open interval of length \(\ell\ge2s\) contains at least
\(\ell/(2s)\) integers in any given residue class modulo \(s\).
For odd \(n\) the two interval lengths are \(n/12,n/4\), and the steps
are \(4L,8L\). For even \(n\) the lengths are \(n/16,n/12\), and the
steps are \(4L,2L\). The two choices are independent and the ordered
support recovers them. Both products of lower bounds are \(n^2/(6144L^2)\).
\end{proof}
A reducible candidate from Theorem~\ref{thm:parity} has at least three
distinct factors, hence a factor of degree between \(B+1\) and
\(\lfloor n/3\rfloor\). Taking \(B(B+1)\le\varepsilon\log_2 n\) as above
gives at least \(n^{2-\varepsilon}/6144\) candidates with the stated
parity and a growing cutoff. The discriminant argument and counting
in this theorem are proved on paper. In particular, this family is
disjoint from the family in Theorem 1, where \(\{0,u+v\}\) and
\(\{u,v\}\) have equal sums.
The two constructions in the proof give the explicit primitive examples
\[
 X^{200}+X^{56}+X^{17}+X^2+1,\qquad
 X^{521}+X^{416}+X^{282}+X^8+1.
\]
Lean certifies both orders and both Sidon criteria in
\path{GF2Sparse/ParityExamples.lean}.
The script \path{scripts/check_parity_sieve.py} independently verifies
one sample in each degree 129--384 at \(B=1\), 385--544 at \(B=2\), and
2689--2696 at \(B=3\) by distinct-degree factorization.

\section{Restrictions on the exponents}
\begin{lemma}\label{lem:compression}
Let \(m\ge1\), \(P=\mathop{\rm lcm}_{1\le d\le m}(2^d-1)\).
Suppose moving each exponent of \(f\) by an integer multiple of \(P\) gives a
polynomial \(g\) of degree \(m\) with constant coefficient 1. Terms with the
same new exponent are added over \(\F\).
Then \(f\) has an irreducible factor of degree at most \(m\).
\end{lemma}
\begin{proof}
Choose an irreducible divisor \(p\) of \(g\) and a root \(\alpha\) of \(p\).
It is nonzero and its order divides \(P\).
Then \(f(\alpha)=g(\alpha)=0\), so \(p\mid f\).
\end{proof}

\paragraph{Width of the three interior terms.}
Suppose \(n>B\) is divisible by \(L\) from Theorem~\ref{thm:sieve}. If
\(0<a<b<c<n\) and \(c-a\le B\), take a root \(\alpha\) of an irreducible
factor of \(1+X^{b-a}+X^{c-a}\). Then \(\alpha^n=1\) and
\[
 (X^n+1+X^a+X^b+X^c)(\alpha)=0.
\]
This gives a proper factor regardless of how far the three interior terms are
shifted. In degrees divisible by \(8L\), trinomials are also ruled out by Swan's
theorem. At degree 24 the least possible interior width is exactly three,
attained by the example in Theorem 1.
\href{https://msp.org/pjm/1962/12-3/pjm-v12-n3-p27-p.pdf}{Swan, \emph{Factorization of polynomials over finite fields}}.

\begin{theorem}\label{thm:bands}
Fix a finite set \(S\subset[0,1]\) containing 0 and 1 and a residue class
\(r\bmod D\), where \(D\ge1\). There are a constant \(c_{S,D}>0\)
and infinitely many degrees \(n\equiv r\bmod D\) such that every irreducible
degree-\(n\) binary polynomial has a nonzero coefficient at an exponent \(e\)
satisfying
\[
                 \min_{s\in S}|e-sn|>c_{S,D}\sqrt{\log n}.
\]
For rational \(S\) and a fixed error allowance \(B\), failure of the restriction
\(\min_{s\in S}|e-sn|\le B\) for all terms occurs on an entire arithmetic progression.
\end{theorem}
\begin{proof}
Let \(B\) range over the positive integers. Write \(k=|S\setminus\{0,1\}|\).
Let \(\delta\) be the minimum of 1 and all distances of interior members of \(S\)
to the two endpoints. Let \(m\) be the least integer at least \(2(B+1)/\delta\)
in the chosen residue class, and take \(P\) as in Lemma~\ref{lem:compression}.
There is a multiple \(D\le t\le DP^k\) of \(D\) and integers \(a_s\)
with \(|ts-a_s|\le 1/P\) for every interior \(s\).
Indeed, put the fractional parts of the vectors \((jDs)_s\), \(0\le j\le P^k\),
in the \(P^k\) half-open cubes of side \(1/P\) in the unit cube. Subtract two
indices in the same cube and multiply their difference by \(D\) to get \(t\).
For \(k=0\) take \(t=D\). Put \(a_0=0,\ a_1=t\) and \(n=m+Pt\equiv r\bmod D\).

Assume every nonzero term of a monic degree-\(n\) polynomial \(f\) with
\(f(0)=1\) is within \(B\) of some \(sn\). For each such term choose one such
\(s\), using 0 and 1 at the endpoints. Replace its exponent \(e\) by
\(b=e-Pa_s\). For interior \(s\),
\[
 |b-sm|\le |e-sn|+P|ts-a_s|\le B+1,
\]
so \(0<b<m\). For \(s=0\) the low exponents are unchanged.
For \(s=1\) the high exponents shift down by \(n-m\).
Consequently 0 and \(m\) each occur exactly once after compression.
Lemma~\ref{lem:compression} gives a divisor of degree at most \(m<n\).

Now \(P\ge2^m-1\) and \(P\le\prod_{d=1}^m(2^d-1)\le2^{m(m+1)/2}\).
Thus \(n\) is unbounded as \(B\) grows and \(\log n=O_{S,D}(B^2)\).
An irreducible polynomial in these degrees must have an exponent farther than
\(B\) from \(nS\), proving the stated bound. For rational \(S\), take \(t\)
to be any positive multiple of \(D\) times a common denominator and \(a_s=ts\) exactly.
\end{proof}
The algebraic compression lemma and the rational-grid obstruction are formalized
in Lean. The simultaneous approximation argument above gives the real-slope extension.

\section{A larger order bound at powers of two}
This application of Gao's iterate method gives information about the order of
the very same candidates that pass the small-factor tests.
\href{https://www.math.clemson.edu/~sgao/papers/Gao99.pdf}{Gao, \emph{Elements of provable high orders in finite fields}}.

\begin{theorem}
Fix \(B\ge1\) and \(L\) as in Theorem~\ref{thm:sieve}.
Let \(n=2^t>L+2\), \(g=(X^L+1)(X^2+1)\), \(d=L+2\), and \(f=X^n+g\).
For every positive integer \(s\) with \(s d^s\le n\),
\[
   \ord_{\F[X]/f}(X)\ \ge\ d^{s(s+1)/2}.
\]
The same bound holds for the order modulo any irreducible factor of degree
\(m\) if \(s d^s\le m\).
\end{theorem}
\begin{proof}
Write \(G_0=X\) and \(G_j=g^{\circ j}\). In the quotient, \(X^n=g(X)\).
Raising to the \(n\)-th power commutes with evaluating any polynomial over
\(\F\), which gives \(X^{n^j}=G_j(X)\) by induction.

Each \(G_j\) has a root that is a root of none of \(G_0,\ldots,G_{j-1}\).
Work in an algebraic closure. For \(j=0,1\) use 0 and 1.
For \(j=2\) choose a nonzero root of \(g+1\). Such a root exists because
\(g+1=X^{L+2}+X^L+X^2\) has distinct exponents and becomes a nonconstant
polynomial with constant coefficient 1 on removing the lowest power of \(X\).
This root maps to 1 and then to 0. Each further preimage of a point outside
\(\{0,1\}\) stays outside that set, since \(g(0)=1\) and \(g(1)=0\).
Taking one preimage at every step increases the first hitting time of 0 by one.
It follows that products of the \(G_j\) with different exponent vectors are
distinct. Cancel common powers and evaluate at such a root for the largest
remaining index.

Take products \(\prod_{j=0}^{s-1}G_j^{e_j}\) with \(0\le e_j<d^{s-j}\).
Their degrees are strictly less than \(s d^s\), since \(\deg G_j=d^j\).
All \(d^{s(s+1)/2}\) products remain distinct modulo \(f\), and all are
powers of \(X\). The same argument works with any stated factor in place of \(f\).
\end{proof}
For fixed \(B\) this gives \(\log\ord(X)\gg_B(\log n)^2\) along powers of two,
by taking \(s=\lfloor\tfrac12\log_d n\rfloor\) for sufficiently large \(n\).
The explicit order inequality, including the version for factors, is formalized
in Lean. The quotient ring may have zero divisors.
One can also let the cutoff grow with the degree. For instance, take
\(B=\lfloor(\log n)^{1/3}\rfloor\) at sufficiently large powers of two.
The product bound for \(L\) gives \(\log d=O(B^2)\). With the same choice of \(s\),
\(s d^s\le n\) and \(\log\ord(X)\gg(\log n)^2/\log d\gg(\log n)^{4/3}\).

\section{Exact primitivity under substitution}
\begin{lemma}\label{lem:sampling}
Let \(F,T\) be monic binary polynomials of the same positive degree \(n\),
and \(m\ge1\). If \(T(X)\mid F(X^m)\), then
\[
 F\text{ is primitive}\quad\Longleftrightarrow\quad
 T\text{ is primitive and }\gcd(m,2^n-1)=1 .
\]
\end{lemma}
\begin{proof}
Write \(N=2^n-1\). A root of order \(N\) has degree at least \(n\),
since a smaller field has fewer than \(N\) nonzero elements.
If \(F\) is primitive, take any root \(\beta\) of \(T\).
Its degree is at most \(n\) and \(\beta^m\) has order \(N\).
Thus \(\beta\) has order exactly \(N\), \(\gcd(m,N)=1\), and \(T\) is
its minimal polynomial. Conversely, if \(\beta\) is a primitive root
of \(T\) and \(\gcd(m,N)=1\), then \(\beta^m\) is a root of \(F\)
of order \(N\).
\end{proof}

\begin{theorem}\label{thm:progression}
Let \(\mathcal P_n\) consist of the degree-\(n\) primitive binary
pentanomials whose support contains four terms in arithmetic progression.
For even \(n\), \(\mathcal P_n\) is empty. For odd \(n\ge5\), cubing
roots gives a bijection from the degree-\(n\) primitive trinomials
onto \(\mathcal P_n\).
\end{theorem}
\begin{proof}
For \(u,m>0\) put
\[
 T=X^a+X^t(1+X^{mu}),\qquad F=X^a+X^t(1+X^u)^m .
\]
The difference-of-powers identity gives \(T\mid F(X^m)\). If
\(a\ne t+mu\) both have degree \(\max(a,t+mu)\). For \(m=3\) the
four block exponents of \(F\) are \(t,t+u,t+2u,t+3u\). Every member
of \(\mathcal P_n\) has this form with \(a\) outside the block.
Deleting the two inner terms gives a primitive trinomial by
Lemma~\ref{lem:sampling}, and \(3\nmid 2^n-1\) forces \(n\) odd.

Conversely, take a primitive trinomial with exponents \(0,k,n\).
Two exponents agree modulo 3, since if all three residues occurred the
trinomial would vanish at a nontrivial cube root of unity, of degree 2.
Use the congruent pair as \(t,t+3u\) and the remaining exponent as \(a\).
If \(a\) coincided with an inner exponent of the block, the endpoints
would be \(0,n=3u\) and \(k\) would be \(u\) or \(2u\).
Here \(u>1\), since \(n\ge5\). A polynomial \(g(X^u)\) with
\(\deg g=d>0\) and \(u>1\) cannot be primitive. For any nonzero
root \(\beta\), its order is at most \(u(2^d-1)<2^{ud}-1\),
because \(\beta^u\) has degree at most \(d\).
Thus the five terms are distinct. Odd \(n\) makes the resulting \(F\)
primitive by the lemma. Cubing permutes the elements of order \(2^n-1\)
and commutes with the binary Frobenius, so it also permutes their
conjugacy classes. This proves uniqueness in both directions.
\end{proof}
For example, the primitive trinomial \(X^5+X^2+1\) gives
\(X^5+X^4+X^3+X^2+1\).
More generally \(m=1+2^j\) with \(j>0\) gives a four-term block
at offsets \(0,u,2^j u,(2^j+1)u\). The same lemma shows that any
genuine primitive pentanomial of this shape requires a primitive
trinomial in the same degree. Kashiwagi and Uchimura introduced the
forward sampling construction.
\href{https://doi.org/10.9746/sicetr1965.18.747}{Kashiwagi and Uchimura,
\emph{A Simple Method for Obtaining Primitive Pentanomials Over GF(2)}}.

\begin{theorem}\label{thm:dilation}
Let \(f\) be irreducible of degree \(n\) with constant coefficient 1,
let \(\alpha\) be a root, and let \(m\mid N=2^n-1\).
Every irreducible factor of \(f(X^m)\) has degree
\(n\ord(\alpha^{N/m})\), and all factors are distinct. In particular,
putting \(h=\ord(\alpha)\), the full substitution \(f(X^N)\) has
exactly \(N/h\) factors of degree \(nh\).
For any binary polynomial \(f\) of positive degree \(n\) with
\(f(0)=1\), primitivity of \(f\) is equivalent to irreducibility of
\(f(X^{2^n-1})\).
\end{theorem}
\begin{proof}
Take a root \(\beta\) of \(f(X^m)\) and write \(\alpha=\beta^m\).
Its conjugacy class has the same order data as any chosen root of \(f\).
The \(n\)-fold Frobenius sends \(\beta\) to
\(\beta\alpha^{N/m}\) and fixes \(\alpha\). A Frobenius power fixing
\(\beta\) must also fix \(\alpha\), so its exponent must be a multiple
of \(n\). The least positive such exponent is therefore
\(n\ord(\alpha^{N/m})\). The polynomial \(f(X^m)\) divides
\(X^{mN}-1\), which has simple roots since \(mN\) is odd.
For \(m=N\), comparison with the total degree \(nN\) gives the factor
count. A nontrivial factorization of \(f\) persists under the
substitution, which covers reducible \(f\) in the equivalence.
\end{proof}
The general subject is classical.
\href{https://doi.org/10.1112/jlms/s1-30.4.480}{Butler,
\emph{The irreducible factors of \(f(x^m)\) over a finite field}}.
The three results in this section, including the bijection, are
formalized in \texttt{GF2Sparse/Dilation.lean} and
\texttt{GF2Sparse/Progressions.lean}. Full substitution preserves
weight and raises the degree to \(n(2^n-1)\).

\section{A norm test for compositions}
Let \(H\in\F[T]\) be primitive of degree \(r\), and let \(h\in\F[X]\)
be monic of degree \(m>0\) with \(h(0)=0\). Write
\[
 f=H(h),\qquad q=2^r,\quad R=q-1,\quad N=q^m-1,\qquad
 D_t=1+q+\cdots+q^{t-1}\quad(t\ge1).
\]
The polynomial \(f\) has degree \(n=rm\). The following test uses
the norm from \(\mathbb F_{q^m}\) to \(\mathbb F_q\).

\begin{theorem}\label{thm:composition-norm}
Every irreducible factor of \(f\) has degree divisible by \(r\).
If \(f\) is irreducible then
\begin{equation}\label{eq:composition-norm}
                         X^{D_m}\equiv h(X)\pmod f .
\end{equation}
Conversely, if \eqref{eq:composition-norm} holds, then \(f\) is
squarefree and each of its irreducible factors has degree \(rt\) with
\(t\mid m\) and \(\gcd(m/t,R)=1\). In particular, if every prime
dividing \(m\) divides \(R\), this one congruence is equivalent to
irreducibility.
\end{theorem}
\begin{proof}
If \(\alpha\) is a root of any irreducible factor, then
\(\beta=h(\alpha)\) is a root of \(H\). Thus
\(\F(\beta)=\mathbb F_q\subseteq\F(\alpha)\), proving the first
assertion, and \(\beta\) has order \(R\).
If \(f\) is irreducible, the extension
\(\mathbb F_q(\alpha)/\mathbb F_q\) has degree \(m\).
Its monic minimal polynomial for \(\alpha\) is \(h(X)-\beta\), of
degree \(m\). The product of the conjugates is its constant term in
characteristic two. Hence
\(\alpha^{D_m}=\prod_{i=0}^{m-1}\alpha^{q^i}=\beta\), as required.

Suppose now that \eqref{eq:composition-norm} holds.
Since \(H(T)\mid T^R-1\), we have \(h^R\equiv1\pmod f\).
Raising the congruence to the \(R\)-th power gives \(f\mid X^N-1\).
This polynomial has simple roots, all in \(\mathbb F_{q^m}\).
Consequently \([\mathbb F_q(\alpha):\mathbb F_q]=t\) divides \(m\).
Write \(m=tk\). The norm \(y=\alpha^{D_t}\) lies in
\(\mathbb F_q^\times\). The \(m\) Frobenius factors in
\(\alpha^{D_m}\) repeat this block of \(t\) factors \(k\) times,
so \(y^k=\beta\). A \(k\)-th power in a cyclic group of order \(R\)
can have order \(R\) only if \(\gcd(k,R)=1\).
Under the stated prime-divisor hypothesis this forces \(k=1\).
Any irreducible factor then has the full degree \(rm\).
\end{proof}

\begin{corollary}\label{cor:composition-order}
For arbitrary \(m\), the composition \(f\) is primitive if and only if
\eqref{eq:composition-norm} holds and
\[
 X^{N/\ell}\not\equiv1\pmod f
 \quad\text{for every prime }\ell\mid N\text{ with }\ell\nmid R .
\]
\end{corollary}
\begin{proof}
Primitivity implies both conditions. Conversely, let \(x\) be the
residue of \(X\) modulo \(f\). Evaluation at \(h(x)\) embeds
\(\F[T]/(H)\) into \(\F[X]/(f)\), since it is a unital map from
a field to a nonzero ring. Hence \(h(x)\) has order \(R\).
The congruence gives \(x^N=1\); for every prime \(\ell\mid R\)
it also gives \(x^{N/\ell}=h(x)^{R/\ell}\ne1\).
Together the checks force \(x\) to have order \(N\).
Its powers then exhaust the \(N\) nonzero elements of the quotient,
so the quotient is a field and \(f\) is primitive.
\end{proof}
Lean verifies Theorem~\ref{thm:composition-norm} and
Corollary~\ref{cor:composition-order} in
\texttt{GF2Sparse/CompositionNorm.lean}.

The test applies to five-term polynomials with two pairs of exponents
related by powers of two. For \(s,i,j,a,b>0\) put \(Q=2^s\) and
\[
 h(X)=\sum_{u=0}^{i-1}X^{aQ^u}+\sum_{v=0}^{j-1}X^{bQ^v}.
\]
Cancellation in \(h^Q+h\) gives
\[
 (T^Q+T+1)\big|_{T=h(X)}=1+X^a+X^{aQ^i}+X^b+X^{bQ^j}.
\]
Four distinct positive exponents give weight five. The outer
polynomials for \(s=1,2\) are primitive of degrees 2 and 4.
Indeed, \(T^4+T+1\) is 1 on \(\mathbb F_4\), hence irreducible,
and differs from the degree-four cyclotomic polynomial for order 5.
For \(s\ge3\) every root of \(T^Q+T+1\) satisfies
\(z^Q=z+1\) and \(z^{Q^2}=z\), so has degree at most \(2s<Q\).
The outer polynomial and its composition with any nonconstant \(h\)
are then reducible. The endpoint identity and this obstruction are
formalized in \texttt{GF2Sparse/FrobeniusPairs.lean}.

Even relative degree \(m\) requires a longer chain for primitivity at
weight five. If \(i=j=1\), then \(h=X^m+X^a\) with \(0<a<m\).
A primitive composition makes \(h(X)-\beta\) a primitive trinomial
over \(\mathbb F_q\), of even base-field degree \(r=2\) or 4.
Vishne's Corollary 5.2 gives \(m=2\). The two outer polynomials
then give \(X^4+X+1\) and
\((X^4+X^3+1)(X^4+X^3+X^2+X+1)\).
With chain lengths 1 and 2, \(h=X^4+X^3+X\) gives the primitive
\(X^{16}+X^{12}+X^3+X+1\).
\href{https://u.math.biu.ac.il/~vishneu/publications/trinomials.pdf}
{Vishne, \emph{Factorization of Trinomials over Galois Fields of Characteristic 2}}.

For \(s=2\), Theorem~\ref{thm:composition-norm} tests irreducibility
in one exponentiation whenever \(m=3^a5^b\) with \(a,b\ge0\).
For example, \(h=X^5+X^4+X\) gives
\(f=X^{20}+X^{16}+X^5+X+1\).
Here \(X^{69905}\equiv h\pmod f\) and
\(2^{20}-1=3\cdot5^2\cdot11\cdot31\cdot41\). The remaining
three powers \(X^{N/11},X^{N/31},X^{N/41}\) have respective
binary coefficient masks \(\mathtt{6DA4},\mathtt{C6A9E},\mathtt{19BC8}\)
in hexadecimal, so the example is primitive.
For general \(m\), the congruence still supplies the full 3-part and
5-part of the order of \(X\) whenever \(f\) is irreducible.
The relative-degree hypothesis matters even for five distinct terms
with exponent gcd 1. Take \(h=X^{46}+X\). Then
\[
 f=h^4+h+1=X^{184}+X^{46}+X^4+X+1.
\]
Modular powering gives \(X^{(2^{184}-1)/15}\equiv h\pmod f\) and
\(X^{2^{92}}\equiv X\pmod f\). Each irreducible factor has degree
dividing 92 and divisible by 4, so its degree is 4 or 92.
Every \(z\in\mathbb F_{16}\) satisfies \(z^{46}=z\), since
\(46=3\cdot15+1\). Thus \(h(z)=0\) and \(f(z)=1\), excluding
degree 4. The polynomial is squarefree, so it has exactly two
degree-92 factors. The examples in degrees 16, 20 and 184 are checked in
\texttt{GF2Sparse/NormExamples.lean}.
Existence of a passing primitive pentanomial in each degree remains open.

\section{Which degrees admit a sparse quadratic lift?}
Substituting \(X^2+X\) doubles the degree. A binary polynomial is of
this form precisely when \(f(X+1)=f(X)\). Indeed, putting \(T=X^2+X\),
one has \(\F[X]=\F[T]\oplus X\F[T]\). The relation \(X^2=T+X\)
gives spanning, and the two summands have different degree parities.
Translation fixes \(T\) and sends \(A(T)+XB(T)\) to
\(A(T)+(X+1)B(T)\). Translation invariants of irreducible
polynomials over finite fields have been studied more generally.
\href{https://arxiv.org/abs/1608.03915}{Reis, \emph{The action of
\(\mathrm{GL}_2(\mathbb F_q)\) on irreducible polynomials over \(\mathbb F_q\)}}.

\begin{theorem}\label{thm:translation}
If a translation-invariant irreducible binary polynomial has at most
five terms and degree \(n\), then
\[
       n\in\{2,4,6,8\}\quad\text{or}\quad n=2^u+4\quad(u\ge3).
\]
For each of the latter degrees, every possible support occurs in one
of the following two families, where \(Q=2^u\):
\[
\begin{aligned}
 &1+X^5+X^{2^j}+X^{Q+1}+X^{Q+4},
                    &&0\le j\le u,\quad j\text{ odd},\\
 &1+(X^Q+X^{2^j})(X^4+X),
                    &&0<j<u,\quad j\ne2.
\end{aligned}
\]
These are necessary conditions. Primitive examples occur in degrees
2, 4, 6, 8, 12 and 20.
\end{theorem}
\begin{proof}
A translation invariant of positive degree has even degree. An
irreducible of such degree has constant term 1 and odd weight.
Write \(Y_i=X^{2^i}\). The binary digits of each exponent specify a
distinct squarefree monomial in the \(Y_i\). Translation replaces
each \(Y_i\) by \(Y_i+1\). The resulting multilinear expressions
can be compared term by term since different subsets give different
exponents of \(X\).

Suppose the largest monomial has \(w\ge2\) variables. In its
translation difference each of its \(w\) faces obtained by deleting
one variable must cancel against a face of a different top monomial.
One other monomial can cover at most one such face. Thus at least
\(w+1\) nonconstant monomials are needed. For \(w=3\) all four
would be triples. Starting with \(abc\), the other three would have
to be \(bcd,ace,abg\). Cancelling the pairs \(cd\) and \(bd\)
forces \(d=e=g\). Each vertex now occurs in three triples, so its
linear coefficient in the translation difference is 1. This rules
out \(w=3\), and hence every exponent has at most two binary digits
equal to 1.

For \(w=2\), form the simple graph whose edges are the quadratic
monomials. Cancellation of every linear coefficient means all vertex
degrees are even. With at most four edges this is either a triangle
or a four-cycle. At weight five the triangle has one additional
linear monomial. Both shapes are invariant: a triangle contributes
constant 1 to the translation difference, which the linear monomial
cancels; a four-cycle contributes 0. The other weight-five shape is
four linear monomials. A trinomial has only two nonconstant terms,
so they are linear. Irreducibility forces one to be \(X\), since
otherwise it is a square. The resulting \(X^{2^s}+X+1\) is
irreducible only in degrees 2 and 4 by the preceding section.

Write \(f(X)=H(T)\) and \(n=2m\). Irreducibility forces \(H\) to
be irreducible. If \(\beta\) is a root of \(H\) in
\(K=\mathbb F_{2^m}\), then \(Z^2+Z+\beta\) must be irreducible
over \(K\), or \(f\) would have a root in \(K\). The image of
\(z\mapsto z^2+z\) in \(K\) is the trace-zero hyperplane: its
kernel has size two, and the absolute trace vanishes on the image
and is nonzero as a map. Hence \(\operatorname{Tr}(\beta)=1\), so
the coefficient of degree \(m-1\) in \(H\) is 1. The two highest
coefficients below the leading term of \(f\) are therefore
\[
       [X^{n-1}]f=m\bmod2,\qquad
       [X^{n-2}]f=1+\binom m2\bmod2.
\]
If \(8\mid n\), the term \(X^{n-2}\) must occur. For \(n>8\)
divisible by eight with at most two set bits, \(n-2\) has at least
three, which is impossible. For \(n>8\) with \(n\equiv2\bmod8\),
one must have \(n=Q+2\), \(Q=2^u\). Both \(Q+1\) and \(Q\) must
occur. The graph must be a triangle with edges of exponents
\(Q+2,Q+1,3\) and linear monomial \(X^Q\). But then
\[
 f=X^{Q+2}+X^{Q+1}+X^Q+X^3+1=(X^2+X+1)(X^Q+X+1).
\]
The only possible degree congruent to 6 modulo 8 is 6 itself.

The remaining degrees above 8 are \(Q+4\). The coefficient of
\(X^{Q+2}\) is zero. The vertex \(u\) must therefore have neighbors
2 and 0: its largest neighbor is 2 and 1 is forbidden. A triangle
gives the first displayed shape. Evaluating at a nontrivial cube
root of unity rules out even \(j\), since \(X^4+X\) vanishes
there and \(1+X^5+X^{2^j}=0\) for even \(j\).
A four-cycle gives the second shape, with \(j\) opposite \(u\).
The small primitive examples have respective supports
\(\{0,1,2\},\{0,1,4\},\{0,2,3,5,6\},\{0,3,5,6,8\},
\{0,5,8,9,12\},\{0,3,6,17,20\}\).
\end{proof}

\begin{theorem}\label{thm:lift-order}
Let \(H\in\F[X]\) be irreducible of degree \(m\), put \(q=2^m\), and
let \(t=[X^{m-1}]H\). Then \(H(X^2+X)\) is irreducible if and only
if \(t=1\). In that case, for any root \(x\) of the lift,
\[
 H(X^2+X)\text{ is primitive}
 \ \Longleftrightarrow\ H\text{ is primitive and }\ord(x^{q-1})=q+1 .
\]
If \(q+1\) is prime and \(H\) is primitive, the single coefficient
\(t=1\) suffices for the lift to be primitive.
\end{theorem}
\begin{proof}
Writing \(\beta=x^2+x\), the sum of the \(m\) Frobenius conjugates
of \(\beta\) telescopes to \(x^q+x=t\). Since \(\beta\) has degree
\(m\), the degree of \(x\) is either \(m\) or \(2m\), and it is \(m\)
exactly when \(x^q=x\). This proves the irreducibility assertion.
When \(t=1\), \(x^{q+1}=x(x+1)=\beta\). The order of \(x\) divides
\((q-1)(q+1)\), and \(q-1,q+1\) are coprime. Thus \(x\) has full
order precisely when \(x^{q+1}\) and \(x^{q-1}\) have orders
\(q-1\) and \(q+1\), respectively. If \(q+1\) is prime, the latter
order is automatic since \(x^q=x+1\ne x\). The trace and order
criteria are formalized in \texttt{GF2Sparse/QuadraticLift.lean}.
\end{proof}

\begin{corollary}\label{cor:mobius}
Suppose a primitive binary polynomial of degree \(n>3\) and weight
at most five has a root set preserved by a nonidentity binary
Möbius transformation. Either the polynomial or its reciprocal is
translation invariant. In particular, above degree 8 the only possible
degrees are \(n=2^u+4\), and there are at most \(4u\) possible supports
including reciprocals.
\end{corollary}
\begin{proof}
For a primitive root \(\alpha\), the product of any nonempty proper
subset of its conjugates is \(\alpha^{\sum_{i\in S}2^i}\ne1\), since
the exponent lies strictly between 0 and \(2^n-1\).
Inversion pairs \(x\) with \(1/x\), whose product is 1. The two elements
of order three in \(\mathrm{PGL}_2(\F)\) cycle
\(x,1/(x+1),(x+1)/x\), also of product 1. These points are distinct
for a root of degree greater than two. Thus none of these three
transformations preserves the root set. The two remaining
nonidentity maps are \(T(x)=x+1\) and \(x/(x+1)\), the conjugate of
\(T\) by inversion. Taking reciprocals preserves degree, weight and
primitivity, so Theorem~\ref{thm:translation} gives the bound.
\end{proof}

\section{An exterior-power test for sparse companion matrices}
\begin{theorem}\label{thm:exterior}
Let \(f\in\F[X]\) be monic of positive degree \(n\) with \(f(0)=1\).
Let \(C\) be its companion matrix. Then
\[
 f\text{ is primitive}\quad\Longleftrightarrow\quad
 \det(I+\textstyle\bigwedge^k C)=1
 \quad(1\le k\le\lfloor n/2\rfloor).
\]
For each even \(n=2m\ge4\) there is an irreducible degree-\(n\)
polynomial passing every test with \(k<m\) and failing at \(k=m\).
\end{theorem}
\begin{proof}
Over a splitting field \(C\) is triangularizable with diagonal the
roots of \(f\), counted with multiplicity. The diagonal entries of its
\(k\)-th exterior power are the products of \(k\) distinct diagonal
positions. Thus the determinant vanishes exactly when such a product is
1. If \(f\) is reducible, an irreducible factor of degree at most \(n/2\)
has constant term 1, so its roots give a forbidden product.

If \(f\) is irreducible, its conjugates are \(\alpha^{2^i}\) for
\(0\le i<n\). For primitive \(\alpha\), each nonempty proper subset has
exponent strictly between 0 and \(2^n-1\), so all the tests hold.
Conversely, if \(\alpha\) has
order \(e<2^n-1\), the binary digits of \(e\) select a nonempty
proper subset with product 1. Its complement also has product 1,
so one of these subsets has size at most \(\lfloor n/2\rfloor\).

For sharpness, choose an element of order \(e=(2^{2m}-1)/3\)
in \(\mathbb F_{2^{2m}}^\times\). It has degree \(2m\) because
\(e>2^m-1\) and every proper subfield has at most \(2^m\) elements.
The only positive multiples of \(e\) below \(2^{2m}-1\) are \(e\)
and \(2e\). Their binary digits occupy respectively the even and odd
positions, exactly \(m\) each. The middle exterior power therefore
has a two-dimensional fixed space. For a five-term example,
\(X^8+X^7+X^3+X+1\) has root order \(85=(2^8-1)/3\).
Lean verifies the half-degree product test and constructs an irreducible
polynomial whose first failed product test has size \(m\) in every
degree \(2m\ge4\), in \texttt{GF2Sparse/ConjugateProducts.lean}.
\end{proof}

For \(f=X^n+\sum_{i\in A}X^i\) the test has an explicit sparse matrix
on \(k\)-subsets \(S\) of \(\{0,\ldots,n-1\}\). If \(n-1\notin S\),
the column of \(\bigwedge^k C\) has a single entry at \(S+1\).
Otherwise put \(D=(S\setminus\{n-1\})+1\). The entries are at
\(D\cup\{i\}\) for \(i\in A\setminus D\). This follows by wedging
the columns of the companion matrix. A pentanomial has \(|A|=4\),
so every column has at most four entries before adding the identity.
Over \(\F\) the determinant counts directed cycle covers modulo two.
Thus it suffices in each degree to find one support for which all
these explicit matrices have odd cycle-cover count. The script
\texttt{scripts/check\_exterior.py} compares this test with the
prime-divisor order test on all 4,095 monic binary polynomials of
positive degree at most 12 with constant coefficient one.

\section{Passing the first two exterior-power tests in every large degree}
\begin{theorem}\label{thm:reciprocal-sieve}
For each \(B\ge1\) there is a constant \(c_B>0\) such that, in every
sufficiently large degree \(n\), at least \(c_B\varphi(n)\) distinct pentanomials \(F_{n,u,v}\)
are squarefree, have exponent gcd 1, avoid all irreducible factors
of degree at most \(B\), and have no pair of mutually inverse roots.
In particular, every one of them passes the tests at \(k=1,2\) in
Theorem~\ref{thm:exterior}.
\end{theorem}
\begin{proof}
Let \(L=\operatorname{lcm}_{1\le d\le B}(2^d-1)\), and choose a
power of two \(P>L\). For positive odd \(w\) coprime to \(n\) in the
interval
\[
 \frac{n}{2P}<w<\frac{n}{P+L},
 \qquad u=Lw,\qquad v=2n-(2P+L)w,
\]
one has \(0<u<v<u+v<n\). Also \(\gcd(n,u,v)=1\), since
\(\gcd(L,2P+L)=1\). The factor bound follows as in
Theorem~\ref{thm:sieve} because \(L\mid u\).

If \(x\) and \(1/x\) were both roots, multiplying the equation
for \(1/x\) by \(x^{u+v+n}\) and using the equation for \(x\) gives
\(x^{2n}=x^{u+v}\). Hence \(x^{2Pw}=1\), which implies \(x^w=1\)
in characteristic two. This contradicts \(F_{n,u,v}(x)=0\),
since \(w\mid u\). Lean checks the inverse-root argument, the
squarefreeness equivalence below, and the primitive example \(F_{55,21,25}\) in \texttt{GF2Sparse/ReciprocalSieve.lean}.

To control repeated roots, put \(c=P+L\) for even \(n\), and \(c=P\)
for odd \(n\). Define
\[
               H(Y)=Y^{P+L}+Y^{2L}+1 .
\]
We claim \(F_{n,u,v}\) is squarefree exactly when
\(\gcd(H(Y),Y^{n-cw}-1)=1\).
Indeed, write \(\beta=x^w\) and \(F=F_{n,u,v}\). For even \(n\),
\(XF'=X^u+X^v\), so at a nonzero repeated root \(x^{v-u}=1\).
This gives \(x^{n-cw}=1\), since \(v-u=2(n-cw)\), and substitution
in \(F\) gives \(H(\beta)=0\). For odd \(n\),
\(F+XF'=1+X^{u+v}\), with \(u+v=2(n-cw)\). Now \(x^n=\beta^P\)
and \(x^v=\beta^{-L}\). Multiplying \(F(x)\) by \(\beta^L\) again
gives \(H(\beta)\). Conversely, since \(\gcd(w,n-cw)=1\),
raising to \(w\) permutes the roots of \(Y^{n-cw}-1\). A common
root \(\beta\) therefore has a preimage \(x\) in this group, and the
same equations give \(F(x)=F'(x)=0\).

Every root of \(H\) has odd order dividing the fixed integer
\(M=\operatorname{lcm}_{1\le d\le P+L}(2^d-1)\).
Also \(H\) has no root of order dividing \(c\). For \(c=P\) the
only candidate is 1. For \(c=P+L\) substituting \(Y^c=1\) leaves
\(Y^{2L}\ne0\). Thus \(\gcd(M,n-cw)\mid c\) suffices.
For each odd prime \(p\mid M\), choose a nonzero class for \(w\bmod p\)
such that \(p^{v_p(c)+1}\nmid n-cw\). At most one class is forbidden
by this condition, since \(c/p^{v_p(c)}\) is a unit modulo \(p\).
Together with oddness these choices specify a unit class modulo
\(D=2\prod_{p\mid M}p\), all of whose members satisfy the sufficient
condition.

The interval has length \(\delta n\) with
\(\delta=1/(P+L)-1/(2P)>0\). Inclusion--exclusion over primes
dividing \(n\) outside \(D\) counts the coprime members of the class as
\[
 \frac{\delta n}{D}\prod_{\substack{p\mid n\\p\nmid D}}(1-1/p)
       +O(2^{\omega(n)}).
\]
Here the main term is at least \(\delta\varphi(n)/D\), and
\(2^{\omega(n)}=o(\varphi(n))\) by the product formula
\(\varphi(n)/2^{\omega(n)}=\prod_{p^a\parallel n}p^{a-1}(p-1)/2\).
Distinct choices of \(w\) have distinct smallest interior exponents
\(u=Lw\), proving the count.
\end{proof}

\begin{corollary}\label{cor:reciprocal-obstruction}
For fixed \(L,P\) in Theorem~\ref{thm:reciprocal-sieve}, put
\(G(Y)=Y^{2(P+L)}+Y^L+1\). For every admissible \(w\),
\[
 \deg\gcd(F_{n,u,v}(X),X^n-1)=\deg\gcd(G(Y),Y^n-1).
\]
In particular, every finite list of fixed choices of \(L,P\) has an
infinite progression of degrees in which every one of its candidates
is reducible.
\end{corollary}
\begin{proof}
In \(R=\F[X]/(X^n-1)\) one has
\[
          X^{(2P+L)w}F_{n,u,v}(X)=G(X^w),
\]
since \((2P+L)w+v=2n\). The power multiplying \(F_{n,u,v}\) is a
unit, and \(X\mapsto X^w\) is an automorphism of \(R\) because
\(\gcd(w,n)=1\). Hence \(R/(F_{n,u,v})\) and \(R/(G)\) have the same
dimension over \(\F\), giving the displayed equality even when
\(X^n-1\) has repeated roots.

Choose one root of \(G\) for each fixed choice of parameters.
Their orders are finite, since \(G(0)=1\). For every degree \(n\)
divisible by all these orders, the gcd is nonconstant for every
candidate. An irreducible degree-\(n\) candidate dividing the monic
degree-\(n\) polynomial \(X^n-1\) would equal it, whereas
\(F_{n,u,v}(1)=1\). For instance \(L=1,P=2\) gives \(G=Y^6+Y+1\)
of root order 63, ruling out all multiples of 63. Lean checks the
coprimality equivalence and the progression obstruction in
\texttt{GF2Sparse/ReciprocalSieve.lean}.
\end{proof}

\section{An exact third exterior-power test for trinomials}
\begin{theorem}\label{thm:triple-orbits}
Let \(T=X^n+X^a+1\) with \(0<a<n\) and \(\gcd(n,a)=1\). Put
\[
       D=n^2-na+a^2,\qquad G=\gcd(T,X^D-1),\qquad
       \epsilon=\mathbf1_{\{a\bmod3,n\bmod3\}=\{1,2\}} .
\]
The roots of \(T\) contain exactly \(\epsilon\) unordered inverse
pairs and \((\deg G-2\epsilon)/3\) unordered triples of distinct roots
with product 1. These triples are pairwise disjoint. In particular
\(G=1\) is the exact test for passing the first three exterior-power
conditions whenever \(n\ge6\).
\end{theorem}
\begin{proof}
The trinomial is squarefree. If just one of \(a,n\) is odd its
derivative is a monomial, and if both are odd then \(T+XT'=1\).
Suppose \(x,y,z\) are roots with \(xyz=1\). Put \(u=x^a,v=y^a,w=z^a\).
Since \(x^n=1+u\), we have \(uvw=(1+u)(1+v)(1+w)=1\). Consequently
\[
             \bigl(v(1+u)+1\bigr)\bigl(w(1+u)+1\bigr)=0 .
\]
Indeed the left side is the sum of \((1+u)(1+v)(1+w)+1\) and
\((1+u)(uvw+1)\). After exchanging \(y,z\) if necessary,
\(y^a=x^{-n}\) and \(y^n=1+y^a=x^{a-n}\).
Computing \(y^{an}\) in two ways gives \(x^D=1\).

Conversely \(\gcd(a,D)=1\). Choose \(j\equiv-na^{-1}\pmod D\).
Then \(jn\equiv a-n\) and \(j^2+j+1\equiv0\pmod D\).
For a root \(x\) of \(G\), the element \(y=x^j\) satisfies
\(y^a=x^{-n}\) and \(y^n=x^{a-n}=1+x^{-n}\), hence is again a root
of \(G\). Since \(j^3\equiv1\pmod D\), this operation partitions the
roots of \(G\) into orbits of size one or three. Each three-cycle has
product 1, and the preceding argument determines that cycle uniquely
from any of its members because powering by \(a\) permutes the
\(D\)-th roots of unity. A fixed point has \(x^3=1\), and \(T(1)=1\).

If \(x,1/x\) are both roots of \(T\), then
\((1+x^a)(1+x^{-a})=1\), so \(x^a\) is a nontrivial cube root of
unity. The same holds for \(x^n=1+x^a\), and \(\gcd(a,n)=1\)
forces \(x^3=1\). The two nontrivial cube roots occur precisely
when \(\epsilon=1\). In this case \(3\mid D\) and \(j\equiv1\bmod3\),
so they give exactly the two fixed points. This proves the counts.
\end{proof}
For instance \(X^{16}+X^3+1\) has \(D=217\), and the gcd has degree
three. Its one product-one triple is the root set of \(X^3+X^2+1\).
The script \texttt{scripts/check\_trinomial\_triples.py} compares both
counts against the nullities of the actual exterior-power matrices.

\begin{corollary}\label{cor:cubic-triples}
For every \(Q=2^s\), the trinomial \(X^{Q+3}+X^Q+1\) has no
product-one triple of roots, including repetitions. The reciprocal
trinomial \(X^{Q+3}+X^3+1\) and odd substitutions of either polynomial
also pass the first three product tests.
\end{corollary}
\begin{proof}
For a root \(x\), write \(u=x^3\). Then \(x^Q(u+1)=1\), giving
\(u^Q(u+1)^3=1\). If \(x,y,z\) were a product-one triple, their
cubes \(u,v,w\) and their translates by 1 would each have product 1.
As in Theorem~\ref{thm:triple-orbits}, we may interchange \(v,w\) so
that \(v(u+1)=1\). Then \((v+1)(u+1)=u\), and hence
\[
 v^Q u^3=(u+1)^3=v^Q\bigl(1+(u+1)^3\bigr).
\]
The first equality follows from \(v^Q(v+1)^3=1\).
For the second multiply \(v^Q(u^Q+1)=1\) by \((u+1)^3\).
Canceling \(v^Q\) gives \(u^3=1+(u+1)^3\), or \(u(u+1)=0\),
contrary to \(u^Q(u+1)^3=1\).
Inversion and powering preserve products. A root at 1 gives a
triple, and an inverse pair \(r,r^{-1}\) gives the triple
\(r,r,r^{-2}\). The trinomials are squarefree since their positive
exponents are relatively prime, and odd substitution preserves simple
nonzero roots. The triple exclusion and its preservation under
substitution are proved in \texttt{GF2Sparse/TripleRoots.lean}.
\end{proof}

\section{Five terms passing three tests in every multiple of eight}
Write \(h(X)=(1+X^2)(1+X^3)\).
\begin{theorem}\label{thm:dyadic-triples}
For \(Q=2^s\), \(s\ge3\), let \(F_Q=X^Q+h\).
No three roots of \(F_Q\) have product 1, even if repetitions are allowed.
For every odd \(m\ge1\), \(F_Q(X^m)\) is a squarefree pentanomial
of degree \(mQ\) passing the first three tests of Theorem~\ref{thm:exterior}.
\end{theorem}
\begin{proof}
Suppose \(x,y,z\) are roots with \(xyz=1\).
Put \(u=x+y+z\), \(v=xy+yz+zx\), so these are the roots of
\(H=X^3+uX^2+vX+1\). Let \(A,B,C\) be the elementary symmetric
functions of \(h(x),h(y),h(z)\) in that order, and put \(\Phi=(A,B)\),
\(E=C+1\). These are polynomials in \(u,v\) over \(\F\).
Since \(h(x)=x^Q\) and likewise for \(y,z\), we have
\(\Phi(u,v)=(u^Q,v^Q)\) and \(E(u,v)=0\). Hence
\[
              E(u,v)=(E\circ\Phi)(u,v)=(E\circ\Phi^2)(u,v)=0 .
\]
We show by an exact polynomial calculation that these equations are inconsistent.

Multiplication by \(X\) on the basis \(1,X,X^2\) modulo \(H\) sends a
column \((a,b,c)\) to \(S(a,b,c)=(c,a+vc,b+uc)\). Set
\(w=(I+S^2)(I+S^3)(1,0,0)^{\mathsf T}\).
The matrix with columns \(w,Sw,S^2w\) represents multiplication by \(h\).
Its trace, sum of principal two-by-two minors and determinant are
\(A,B,C\), also for repeated roots, by triangularizing multiplication by \(X\).
In particular, putting \(t=u+v\) gives \(C=t^3(t^2+uv)\).
Since \(C=1\), we have \(t\ne0\). With \(b=u/t\), the first equation becomes
\[
              u=tb,\qquad v=t(b+1),\qquad b^2+b=1+t^{-5}.
\]
Reduce the other two equations in
\(\F[t,t^{-1}][b]/(b^2+b+1+t^{-5})\).
The norm of \(a+db\) is \(a^2+ad+(1+t^{-5})d^2\), its product with
\(a+d(b+1)\). A zero at the alleged solution gives a zero of its norm.
Normalize each norm to lowest exponent zero. For
\(E\circ\Phi\) and \(E\circ\Phi^2\) the normalizing factors are
\(t^{57}\) and \(t^{321}\), and the resulting polynomials \(N_1,N_2\)
have degrees 95 and 535. Encoding the coefficient of \(t^i\) as bit \(i\),
the matrix operations above give the hexadecimal values
\[
\begin{aligned}
 N_1 &= \mathtt{ceea1726aa3d6e9d01a0b261},\\
 N_2\bmod N_1 &= \mathtt{5126499684e6d15d7a2a7ed4}.
\end{aligned}
\]
The Euclidean algorithm gives \(\gcd(N_1,N_2)=1\). Thus there is no common
nonzero \(t\). The complete calculation from the matrix is reproduced by
\texttt{scripts/check\_dyadic\_triples.py}.

The derivative \(F_Q'=X^2(X+1)^2\) has no root in common with \(F_Q\),
since \(F_Q(0)=F_Q(1)=1\). Odd substitution preserves simple roots,
by the chain rule and the nonzero constant term. It also preserves the
exclusion of triples, since raising a product-one triple to its \(m\)-th
powers would give one for \(F_Q\). Root 1 is already excluded. An inverse
pair \(r,r^{-1}\) would give the forbidden triple \(r,r,r^{-2}\),
since squaring preserves binary roots. This proves the three tests.
\end{proof}
The degree \(Q\) supports have exponent gcd 1. The examples \(F_8\)
and \(F_{16}\) are primitive. Lean verifies the Bézout identity in
\texttt{GF2Sparse/DyadicCertificate.lean} and proves the theorem, the
odd dilations and both primitive examples in \texttt{GF2Sparse/DyadicTriples.lean}.
If \(m>1\), a root of \(F_Q(X^m)\)
has order at most \(m(2^Q-1)<2^{mQ}-1\), since its \(m\)-th power
has degree at most \(Q\). The substitutions therefore settle the three
tests even in degrees where further tests necessarily fail.

\paragraph{Degrees near powers of two.}
Put \(F_n=X^n+X^5+X^3+X^2+1\) for \(n>5\).
\begin{theorem}\label{thm:shifted-triples}
Suppose \(n=2^s+d\) with
\[
 d\in\{-3,-2,-1,0,1,3,4,6,7,8,9\}.
\]
Then \(F_n\) passes the first three exterior-power tests exactly when
\(n\not\equiv0,5\pmod7\). In particular it passes in every degree
\(2^s+7\) and \(2^s+9\). All these pentanomials have exponent gcd 1.
\end{theorem}
\begin{proof}
Replace \(h(X)\) in the cubic calculation by \(h(X)/X^d\).
A root of \(F_n\) is taken to its \(2^s\)-th power by this map.
When \(xyz=1\), the product of the three images is still
\(t^3(t^2+uv)\). Thus the same parameterization
\(u=tb,\ v=t(b+1),\ b^2+b=1+t^{-5}\) applies.
For the new image map \(\Phi_d\), compute the normalized norms of
\(E\circ\Phi_d\) and \(E\circ\Phi_d^2\) as above.
Every possible \(t\) is a root of their gcd \(G_d\).

There is a short exact test for the remaining possibilities even when
\(G_d\ne1\). Write \(a=X^3+tX+1\) and \(w=X^2+X\). The cubic
\(H\) is \(a+twb\) and its norm multiplied by \(t^3\) is
\[
           J(t,X)=t^3(a^2+twa+t^2w^2)+w^2.
\]
For each irreducible factor \(g\) of \(G_d\), form
\(\operatorname{Res}_t(g,J)\). Every coordinate of a product-one
triple belongs to the roots of these resultants. Among their distinct
irreducible factors \(p(X)\), retain exactly those dividing \(F_n\).
For a factor of degree \(e\) this test depends only on \(s\bmod e\),
since Frobenius to the \(e\)-th power fixes its roots.
The product of retained factors is squarefree. Testing its first three
exterior powers detects any triple of distinct roots, any inverse
pair, and root 1. This also detects repeated triples, since
\(x^2y=1\) gives an inverse pair or root 1. Conversely an inverse
pair gives a repeated triple; its two roots are retained because
the resultants are binary and their root sets are closed under
squaring and square roots. Thus the finite calculation is exact.

The table gives \(G_d\) in hexadecimal and the passing values of
\(s\bmod3\). The calculation of each gcd and resultant, followed by the
exterior-power tests, is in \texttt{scripts/check\_shifted\_triples.py}.
\[
\begin{array}{r|rrrrrrrrrrr}
d&-3&-2&-1&0&1&3&4&6&7&8&9\\ \hline
G_d &5&11&5&1&5&14115&39&5&1&541&41\\
s\bmod3&1,2&0,2&1,2&0,1,2&0,1&0&1,2&1,2&0,1,2&0,1&0,1,2
\end{array}
\]
These are precisely the indicated congruences on \(n\).
The two excluded classes have cubic factors. Modulo \(X^3+X+1\)
one has \(h=1\), and modulo \(X^3+X^2+1\) one has \(h=X^5\).
The nonzero roots of either cubic have order 7.
\end{proof}
The same script gives exact periods 45 and 24 in \(s\) for the
offsets 2 and 5, respectively. The degree \(265=2^8+9\) member
\(X^{265}+X^5+X^3+X^2+1\) is primitive.

\section{Counting primitive witnesses}
For any set \(\mathcal A\) of monic degree-\(n\) candidates, put \(N=2^n-1\).
An exact expression for the number of primitive witnesses is
\[
 n\,\#\{f\in\mathcal A:f\text{ is primitive}\}
  =\sum_{e\mid N}\mu(e)\sum_{f\in\mathcal A}
      \deg\gcd(f,X^{N/e}-1).
\]
Indeed, \(X^N-1\) has simple roots, precisely the nonzero elements of
\(\mathbb F_{2^n}\). A root of order \(h\mid N\) has total coefficient
\(\sum_{e\mid N/h}\mu(e)\), which is 1 when \(h=N\) and 0 otherwise.
A root of order \(N\) has degree \(n\), so its minimal polynomial accounts
for all \(n\) roots of that candidate. For the two-parameter family in
Theorem 1, strict positivity of the right-hand side in every degree
\(n\ge5\) would suffice. The local sieve and the order bounds leave this
positivity question open.

\paragraph{Counting a whole family with one product.}
Let \(P=\prod_{f\in\mathcal A}f\) for a set of distinct degree-\(n\) candidates.
Start with \(G=\gcd(P,X^N-1)\). For each distinct prime \(p\mid N\) replace
\[
                   G\ \longleftarrow\ G/\gcd(G,X^{N/p}-1).
\]
The initial \(G\) is squarefree. Each division removes precisely the
remaining roots whose orders divide \(N/p\). The final polynomial is
the product of the primitive members of \(\mathcal A\), since each
surviving root has degree \(n\). Its degree divided by \(n\) therefore
counts the witnesses, with no need to factor the individual candidates.

This calculation detects a small exception for the Sidon support restriction.
In degree 12 there are 14 normalized five-term Sidon supports.
The initial gcd has degree 64, and the very first prime, \(p=3\), removes
it completely. Thus none is primitive. In degrees 11 and 13 the same
calculation gives respectively 2 primitive members out of 4, and 14 out of 28.
The command \texttt{python3 scripts/check\_parity\_sieve.py --sidon-count 12}
reproduces the degree-12 count. The unrestricted degree-12 claim already
has a Lean witness in Theorem 1. A separate Lean proof of the restriction
tests the exponents 4095 and 1365 on every admissible support in
\texttt{GF2Sparse/ParityExamples.lean}.

\end{document}
